\begin{finding}[\fid{F22}: constant nearby-cycle purity does not detect collision order]\label{f:collision}
Let $R=k[[t]]$ and, for $m\ge1$, set
\[
 C_m=\{(u,w)\in R^2:u\equiv w\pmod{t^m}\}
     \simeq R[x]/(x(x-t^m)).
\]
For $m\ge n\ge1$, $C_m\subset C_n$ induces a finite universal
homeomorphism $\Spec C_n\to\Spec C_m$ over $\Spec R$, although
\[
 \operatorname{length}_R(R^2/C_m)=m,\qquad
 \operatorname{cond}(C_m,R^2)=t^mR^2.
\]
Their corresponding ordinary \'etale coefficient categories and
cohomology are equivalent. With $k=\F_q$ and unshifted constant
$\Q_\ell$ coefficients, the unique special-point nearby and vanishing
cycles are
\[
 R\Psi\Q_\ell=\Q_\ell^2[0],\qquad
 R\Phi\Q_\ell=\Q_\ell[0],
\]
with trivial inertia and Frobenius. These constant data, including every
finite-extension trace and tensor power, are independent of $m$.
The completed center of \fid{F19} in good characteristic is $C_2$.
\end{finding}
\begin{proof}
Map $x$ to $(0,t^m)$ to obtain the displayed description. For $m\ge n$,
$C_n$ is a finite $C_m$-module since it is already generated by two
elements over the subring $R$. After inverting $t$, the map is the
identity on $k((t))^2$. Each special fiber has one point, with residue
field $k$, and the map induces the identity on that field. It is
therefore integral, surjective and radicial, hence a universal
homeomorphism. The difference map
$(u,w)\mapsto w-u\bmod t^m$ has kernel $C_m$ and gives
$R^2/C_m\simeq R/(t^m)$. The integral closure is $R^2$; a pair multiplies
every element of $R^2$ back into $C_m$ precisely when both coordinates
are divisible by $t^m$. This proves the coherent formulas.

Topological invariance of the \'etale site gives the claimed equivalence,
compatible with the base trait; precise locators are Stacks
Theorem~59.45.2, tag~04DZ, and Proposition~59.45.4, tag~03SI,
the former citing EGA~IV, 18.1.2 \cite{StacksNil}.
This is an equivalence of coefficient-sheaf topoi, not an identification
of coherent structure sheaves as ringed objects.

After strict henselization at the special point, the geometric generic
fiber is two points. Indeed the finite local algebra over the strictly
henselian trait is strictly henselian, and its two branches are already
split over the fraction field. The defining nearby-cycle stalk is
therefore the cohomology of two geometric points: $\Q_\ell^2$ in degree
zero, with trivial inertia and Frobenius. The specialization map from
$i^*\Q_\ell=\Q_\ell[0]$ is diagonal. The defining triangle
\[
 i^*\Q_\ell\longrightarrow R\Psi\Q_\ell
 \longrightarrow R\Phi\Q_\ell\longrightarrow i^*\Q_\ell[1]
\]
thus gives the stated cokernel in degree zero and no other groups.
This also follows by finite base change, Stacks Lemma~59.91.9,
tag~0A4E. Every object and map just described is independent of $m$.
The traces of tensor powers $r$ of $R\Psi$ and $R\Phi$ are respectively
$2^r$ and $1$ over every finite extension.

Finally \eqref{eq:r3collision} says that the actual root difference is
a unit times $(v-1)^2$. Translate $z$ by a root and rescale the translated
coordinate by that unit after completion. The relation becomes
$x(x-t^2)=0$, where $t=v-1$. This identifies the center with $C_2$.
\end{proof}
This is an information limit for the specified constant coefficient
families. A nilpotent algebra action supplied separately, or an
$m$-dependent coefficient construction, is additional information and
is not ruled out. A sign criterion applied only to these common traces
can establish their correct weight-zero purity but cannot determine
$m$. The next finding exhibits coherent information that does retain it.

